% ==============================================================================
% secoes/04_resultados.tex - SEÇÃO IV: RESULTADOS
% Tema: Federico Faggin e o Primeiro Microprocessador Monolítico (Intel 4004)
% ==============================================================================

\section{Resultados}
\label{sec:resultados}

A viabilização física do Intel 4004 em uma única pastilha de silício resultou da convergência entre a tecnologia de porta autoalinhada de silício (SGT), técnicas de \textit{layout} compacto e circuitos dinâmicos de carga. Ao solucionar os entraves eletrostáticos e geométricos da lógica aleatória, Federico Faggin obteve métricas de velocidade e densidade que habilitaram a produção comercial do microprocessador e alteraram a organização da indústria eletrônica. Esta seção reúne os resultados quantitativos alcançados no dispositivo e no circuito, seguidos da análise de seus impactos industriais, econômicos e sociais.

\subsection{Inovações Técnicas de Dispositivo e Circuito}
\label{subsec:contribuicoes_tecnicas}

A aplicação das equações da \cref{subsec:modelagem_fisica_sgt} aos parâmetros nominais de fabricação do Intel 4004 ($t_{ox} \approx \SI{100}{\nano\meter}$ e $\kappa_{ox} \approx 3{,}9$) resulta em uma capacitância específica de óxido $C_{ox} \approx \SI{34,5}{\nano\farad\per\centi\meter\squared}$, equivalente a \SI{0,345}{\femto\farad\per\micro\meter\squared} (\cref{eq:cox}) \cite{faggin1970, sze2006}. Como a porta refratária de polissilício elimina a margem litográfica de alinhamento ($\Delta L_{align} = 0$, \cref{eq:lov_sgt}), a sobreposição longitudinal restringe-se à difusão lateral $L_{ov} \approx \SIrange{0,5}{1,0}{\micro\meter}$, contra \SIrange{3,5}{6,0}{\micro\meter} na porta de alumínio. Pela \cref{eq:cov}, essa redução comprime a capacitância de sobreposição $C_{ov}$ para a faixa de \SIrange{0,17}{0,35}{\femto\farad\per\micro\meter} de largura de canal, ante \SIrange{1,2}{2,1}{\femto\farad\per\micro\meter} no processo metálico. A queda resultante em $C_{gd}$ atenua o carregamento amplificado pelo efeito Miller (\cref{eq:miller}) e reduz o atraso de propagação $\tau_{pd}$ (\cref{eq:tau_pd}), elevando a frequência máxima de relógio em até cinco vezes sob a mesma litografia de \SI{10}{\micro\meter} \cite{faggin1970, faggin1996}.

No plano eletrostático (\cref{eq:vth_expandido,eq:delta_phims}), a combinação do polissilício $p^+$ com lâminas de orientação $\langle 100 \rangle$, cujos parâmetros constam da \cref{tab:comparacao_eletrica}, reduziu o módulo da tensão de limiar $V_{th}$ no Intel 4004 \cite{klein1972patent, sze2006}. Essa redução viabilizou a operação do conjunto MCS-4 com diferença de \SI{15}{\volt} entre trilhos, garantindo compatibilidade direta em nível alto com circuitos bipolares TTL sem conversores externos \cite{intel1971datasheet}. Em paralelo, a captura iônica com $\text{POCl}_3$ em alta temperatura estabilizou $V_{th}$ contra derivas por íons de sódio ($\text{Na}^+$) e reduziu a corrente de fuga reversa das junções em mais de 100 vezes frente ao processo com alumínio \cite{faggin1970}.

A frequência máxima de operação ($f_{max}$) do processador depende das propriedades de transporte dos portadores e das dimensões geométricas conforme:
\begin{equation}
\label{eq:res_fmax}
f_{max} \propto \frac{\mu_p \cdot C_{ox} \cdot (|V_{SS} - V_{DD}| - |V_{th}|)^2}{L^2 \cdot C_L},
\end{equation}
em que $|V_{SS} - V_{DD}| = \SI{15}{\volt}$ representa a tensão efetiva entre os trilhos de alimentação, $L = \SI{10}{\micro\meter}$ é o comprimento do canal e $\mu_p$ ($\approx \SIrange{150}{200}{\centi\meter\squared\per\volt\per\second}$) designa a mobilidade efetiva das lacunas, grandeza que mede a facilidade com que esses portadores positivos atravessam o canal sob campo elétrico, de modo análogo a caminhões carregados em comparação a carros leves representados pelos elétrons \cite{sze2006}. Sob o gerador de relógio bifásico não sobreposto ($\phi_1$ e $\phi_2$), o Intel 4004 atingiu frequência de \SI{740,7}{\kilo\hertz} (especificada nominalmente entre \SI{740}{\kilo\hertz} e \SI{750}{\kilo\hertz}), correspondente a um período elementar de \SI{1,35}{\micro\second} \cite{intel1971datasheet, intel1973mcs4}. Como cada instrução de uma palavra percorre os oito subciclos de A1 a X3 detalhados na \cref{subsec:arquitetura_mcs4_barramento} e na \cref{fig:mcs4_cronograma_subciclos}, o ciclo de instrução resulta em \SI{10,8}{\micro\second}, sustentando uma vazão próxima de \num{92600} instruções por segundo \cite{intel1973mcs4, faggin1996}.

No plano topológico, a densidade bruta de integração ($D_{int}$) alcançada no Intel 4004 para lógica aleatória atingiu:
\begin{equation}
\label{eq:res_densidade_int}
D_{int} = \frac{N_{\text{transistores}}}{\text{Área}} \approx \frac{2300}{\SI{12}{\milli\meter\squared}} \approx \num{192}\ \text{transistores/mm}^2.
\end{equation}
Esse empacotamento em uma pastilha de \SI{3,0}{\milli\meter} por \SI{4,0}{\milli\meter} foi viabilizado pela redução média de 50\% na área das células lógicas propiciada pelos contatos enterrados (\cref{subsec:inovacoes_circuito_leiaute}) e pelo roteamento ortogonal em duas camadas isoladas \cite{vadasz1972patent, faggin2018howwemade, faggin2021}. Sem os contatos enterrados, a área estimada da CPU ultrapassaria \SI{20}{\milli\meter\squared}, comprometendo o rendimento de fabricação (\textit{yield}), isto é, a fração de pastilhas sem defeitos aproveitáveis em cada lâmina de silício, assim como cortar moldes maiores sobre uma folha com pequenas falhas aleatórias aumenta a probabilidade de cada peça conter um defeito e ser descartada \cite{chm2004faggin}.

No circuito das portas lógicas, o desempenho medido nos barramentos internos decorreu da substituição de cargas ativas convencionais pelo estágio com \textit{bootstrapping} dinâmico (\cref{eq:delta_vg_boot,eq:vg_total_boot}). Em uma carga PMOS simples em enriquecimento, a descarga rumo ao trilho negativo $V_{DD}$ estanca em $V_{out,\text{lim}} = V_{DD} + |V_{th}|$ (\cref{eq:vout_max_sem_boot}), próximo a \SI{-8}{\volt}, reduzindo a margem de ruído e retardando o final da transição. Com o acoplamento capacitivo do \textit{bootstrapping}, o potencial de porta $V_G(t)$ mergulha além de $V_{DD}$ em módulo durante toda a descarga, mantendo $|V_{GS}| > |V_{th}|$ e restaurando a excursão integral $V_{out} = V_{DD}$ (\SI{-10}{\volt}) \cite{faggin1972fourchip, faggin1996, weste2010}. A \cref{fig:desempenho_sgt_10um} e a \cref{tab:comparativo_sgt} reúnem o contraste quantitativo obtido sob litografia de \SI{10}{\micro\meter}, enquanto a \cref{tab:inovacoes_faggin} consolida os mecanismos aplicados por Faggin no Intel 4004 e seus resultados de projeto.

\begin{figure}[!t]
\centering
\includegraphics[width=\columnwidth]{desempenho_sgt_10um.pdf}
\caption{Comparação proporcional entre a porta metálica de alumínio e o processo PMOS SGT do Intel 4004 na mesma litografia de \SI{10}{\micro\meter} ($L_{ov}$, $C_{ov}$, $|V_{th}|$, frequência $f_{max}$ e densidade de transistores).}
\label{fig:desempenho_sgt_10um}
\end{figure}

\begin{table*}[!t]
\renewcommand{\arraystretch}{1.2}
\caption{Comparativo Tecnológico Quantitativo entre Porta de Alumínio e Tecnologia SGT}
\label{tab:comparativo_sgt}
\centering
\begin{tabularx}{\textwidth}{>{\raggedright\arraybackslash}p{3.4cm} >{\raggedright\arraybackslash}p{3.8cm} >{\raggedright\arraybackslash}p{4.1cm} >{\raggedright\arraybackslash}X}
\toprule
\textbf{Parâmetro Técnico} & \textbf{Processo PMOS com Porta de Alumínio} & \textbf{Processo PMOS SGT (Intel 4004)} & \textbf{Fundamento Físico e Referência} \\
\midrule
Tolerância de alinhamento ($\Delta L_{align}$) & \SIrange{3}{5}{\micro\meter} & \SI{0}{\micro\meter} (autoalinhado) & Porta de polissilício atua como máscara de difusão (\cref{eq:lov_sgt}) \cite{kerwin1969patent, faggin1970}. \\
Comprimento de sobreposição ($L_{ov}$) & \SIrange{3,5}{6,0}{\micro\meter} & \SIrange{0,5}{1,0}{\micro\meter} & Sobreposição restrita à difusão lateral sob o óxido (\cref{eq:lov_sgt}) \cite{faggin1970}. \\
Capacitância de sobreposição ($C_{ov}$) & \SIrange{1,2}{2,1}{\femto\farad\per\micro\meter} de canal & \SIrange{0,17}{0,35}{\femto\farad\per\micro\meter} de canal & Redução de 60\% a 80\% com $C_{ox} \approx \SI{0,345}{\femto\farad\per\micro\meter\squared}$ (\cref{eq:cov}). \\
Efeito Miller dinâmico ($C_{in,Miller}$) & $(4\text{ a }6) \times C_{gd,metal}$ & $(4\text{ a }6) \times C_{gd,SGT}$ & Amplificação por $(1 + |A_v|)$ sobre $C_{gd}$ três a cinco vezes menor (\cref{eq:miller}) \cite{sze2006}. \\
Tensão de alimentação ($V_{DD}$) & \SIrange{-24}{-28}{\volt} & \SI{15}{\volt} entre trilhos ($V_{SS} - V_{DD}$) & Operação com $V_{SS} = \SI{+5}{\volt}$ e $V_{DD} = \SI{-10}{\volt}$ após redução de $|V_{th}|$ \cite{intel1971datasheet}. \\
Interface com circuitos TTL & Incompatível diretamente, exige conversores & Compatível em nível alto pelo trilho de \SI{+5}{\volt} & Eliminação de tradutores externos com lógica bipolar \cite{intel1971datasheet}. \\
Frequência máxima de relógio & Inferior a \SI{150}{\kilo\hertz} & \SIrange{740}{750}{\kilo\hertz} (\SI{740,7}{\kilo\hertz} nominal) & Ganho de velocidade de até cinco vezes em litografia de \SI{10}{\micro\meter} (\cref{eq:res_fmax}) \cite{faggin1996}. \\
Camadas condutoras de roteamento & Camada única de alumínio & Duas camadas isoladas (polissilício e alumínio) & Roteamento ortogonal duplo elimina túneis difusos \cite{vadasz1972patent, faggin2021}. \\
Densidade bruta de integração & \numrange{90}{100} transistores por \si{\milli\meter\squared} & \num{192} transistores por \si{\milli\meter\squared} (\num{2300} em \SI{12}{\milli\meter\squared}) & Compactação por contatos enterrados e duas camadas condutoras (\cref{eq:res_densidade_int}) \cite{faggin2018howwemade}. \\
Corrente de fuga e estabilidade & Suscetível a desvios iônicos e \textit{spiking} & Fuga mais de 100 vezes menor e limiar estável & Recozimento refratário com captura iônica em $\text{POCl}_3$ \cite{faggin1970}. \\
\bottomrule
\end{tabularx}
\end{table*}

\begin{table*}[!t]
\renewcommand{\arraystretch}{1.2}
\caption{Inovações Técnicas de Faggin no Intel 4004, Mecanismos Físicos e Desdobramentos Industriais}
\label{tab:inovacoes_faggin}
\centering
\begin{tabularx}{\textwidth}{>{\raggedright\arraybackslash}p{4.0cm} >{\raggedright\arraybackslash}p{5.5cm} >{\raggedright\arraybackslash}X}
\toprule
\textbf{Inovação Técnica e Período} & \textbf{Mecanismo Físico ou Arquitetural} & \textbf{Desdobramentos Industriais} \\
\midrule
Porta autoalinhada de silício (SGT, 1968 a 1970) & Porta de polissilício $p^+$ como máscara de difusão térmica e controle de carga interfacial (\cref{eq:lov_sgt,eq:vth_expandido,eq:delta_phims}). & Estreia fabril no Fairchild 3708, adoção na DRAM Intel 1103 e base para a velocidade de operação do 4004 \cite{faggin1970}. \\
Contato enterrado (\textit{Buried Contact}, 1968 a 1970) & Junção ôhmica direta entre polissilício e difusão sem vias para o alumínio, reduzindo a área das células lógicas. & Compactação da pastilha para \SI{3}{\milli\meter} por \SI{4}{\milli\meter}, evitando áreas superiores a \SI{20}{\milli\meter\squared} e preservando o rendimento \cite{vadasz1972patent, faggin2021}. \\
Carga ativa com \textit{bootstrapping} dinâmico (1970 a 1971) & Acoplamento capacitivo em nó flutuante (\cref{eq:delta_vg_boot,eq:vg_total_boot}) mantendo $|V_{GS}| > |V_{th}|$ durante a descarga. & Restauração da excursão plena $V_{out} = V_{DD}$ e ampliação das margens de ruído nos barramentos internos \cite{faggin1996}. \\
Roteamento ortogonal em lógica aleatória (1970 a 1971) & Pistas verticais de polissilício e horizontais de alumínio em duas camadas isoladas, ligadas por vias e contatos enterrados. & Execução do traçado manual da CPU sem ferramentas automáticas, base do projeto VLSI estruturado \cite{mead1980, faggin2021}. \\
Barramento multiplexado em DIP de 16 pinos (1970 a 1971) & Compartilhamento temporal de 4 pinos bidirecionais de dados e controle em oito subciclos (A1 a X3) por instrução. & Viabilização do encapsulamento cerâmico padrão de 16 pinos e operação sincronizada da família MCS-4 \cite{intel1971datasheet, faggin1974patent}. \\
\bottomrule
\end{tabularx}
\end{table*}

\subsection{Impactos na Indústria e na Sociedade}
\label{subsec:impactos_sociedade}

A disponibilidade comercial do Intel 4004 reestruturou o projeto de equipamentos digitais ao substituir a lógica discreta cabeada, formada por circuitos cuja função fica fixada permanentemente nas ligações físicas entre portas lógicas separadas e fios soldados em placa, por arquiteturas programáveis em memória \cite{faggin1972fourchip, cavin2012}. Até o início da década de 1970, desenvolver um controlador exigia interligar dezenas de encapsulamentos TTL SSI e MSI em placas dedicadas. Qualquer ajuste funcional obrigava redesenhar placas e matrizes fabris, gerando elevados custos de engenharia não recorrente (NRE, sigla para o gasto inicial único com projeto, máscaras e ferramental antes da produção em série) e prazos de vários meses. Com o conjunto MCS-4 (CPU 4004, ROM 4001, RAM 4002 e registrador 4003), uma placa substituiu em média noventa chips TTL discretos, transferindo a lógica do produto para o \textit{firmware}, isto é, o programa de baixo nível gravado em memória não volátil para governar o hardware (\cref{fig:impacto_ssi_vs_mcs4}) \cite{faggin1996, intel1988annualreport}. Alterações passaram a ocorrer pela regravação do código em máscaras de ROM ou em memórias EPROM (memórias somente de leitura programáveis e apagáveis por luz ultravioleta, próprias para testes e atualizações rápidas), comprimindo o ciclo de desenvolvimento de meses para semanas e reduzindo custos de manufatura e manutenção \cite{faggin1972fourchip, intel1988annualreport}.

\begin{figure}[!t]
\centering
\includegraphics[width=\columnwidth]{impacto_ssi_vs_mcs4.pdf}
\caption{Transição estrutural promovida pelo conjunto MCS-4: substituição de dezenas de circuitos TTL em lógica discreta cabeada por quatro componentes padronizados cuja função reside no \textit{firmware} armazenado em ROM ou EPROM.}
\label{fig:impacto_ssi_vs_mcs4}
\end{figure}

Essa programabilidade consolidou os sistemas embarcados, computadores dedicados incorporados ao interior de máquinas e equipamentos para executar tarefas específicas de controle. Após estrear na calculadora impressora Busicom 141-PF em 1971, a família MCS-4 avançou sobre segmentos antes atendidos por mecanismos eletromecânicos \cite{chm2007panel, intel1988annualreport}. Na instrumentação biomédica, o analisador clínico de eletroforese (técnica laboratorial que separa proteínas em amostras biológicas mediante campo elétrico) da Helena Laboratories substituiu múltiplas placas dedicadas e uma fonte por um módulo compacto com o Intel 4004, reduzindo o volume físico e o custo de produção do equipamento \cite{intel1988annualreport}. O novo modelo também equipou balanças comerciais com cálculo automático de preço por peso, controladores adaptativos de semáforos, bombas de combustível automatizadas e caixas registradoras. No setor de máquinas recreativas operadas por moedas, o protótipo de pinball \textit{Flicker} (1974), desenvolvido pela Dave Nutting Associates com base no Intel 4004, antecipou as mesas comerciais em estado sólido, termo que designa a eletrônica de semicondutores sem peças mecânicas nem relés. Esse protótipo permanece distinto do posterior modelo comercial \textit{Spirit of 76} da Mirco Games, que adotou o Motorola 6800 \cite{faggin1996, faggin2021, intel1988annualreport}.

Na automação fabril e na robótica, a arquitetura derivada do 4004 viabilizou o controle programável de manipuladores industriais. Em 1974, o robô sueco ASEA IRB 6 estreou como o primeiro manipulador industrial totalmente elétrico comandado por microprocessador, utilizando o Intel 8008, herdeiro direto da metodologia SGT consolidada no 4004. Ao substituir relés eletromecânicos, contadores analógicos e longos chicotes de fiação por controle digital, o sistema elevou o tempo médio entre falhas (MTBF, intervalo médio de operação contínua antes de uma pane reparável) e aumentou a precisão de repetição.

A difusão fabril desses controladores programáveis também alterou a composição do emprego industrial. O estudo de Sar A. Levitan e Clifford M. Johnson, publicado em 1982 na revista \textit{Monthly Labor Review} do Bureau of Labor Statistics, documentou que a adoção de microprocessadores e robôs não provocou desemprego estrutural súbito e generalizado, isto é, perda duradoura de vagas por obsolescência tecnológica de ocupações inteiras, mas desencadeou uma redistribuição gradual de postos de trabalho \cite{bls1982futurework}. A demanda por tarefas manuais repetitivas retraiu-se enquanto cresceram as funções ligadas à programação, manutenção eletrônica e supervisão digital. Essa tendência foi corroborada por estudos econômicos sobre máquinas CNC (máquinas-ferramenta de comando numérico computadorizado, cujos movimentos de corte são guiados por programas digitais em vez de ajustes manuais), cujas instalações registraram maior produtividade acompanhada de maior exigência de qualificação técnica \cite{nber2022cnc}.

A metodologia do Intel 4004 alicerçou ainda a linhagem de microprocessadores de 8 bits conduzida por Faggin na década de 1970, representada na \cref{fig:linhagem_faggin_cpu}. Em 1972, o Intel 8008 empregou o mesmo processo PMOS SGT de \SI{10}{\micro\meter} com \num{3500} transistores e endereçamento de até 16 quilobytes \cite{faggin1996}. Em 1974, Faggin e Masatoshi Shima lideraram o Intel 8080 em tecnologia NMOS SGT de \SI{6}{\micro\meter} a \SI{2}{\mega\hertz}, integrando endereçamento de 64 quilobytes e pilha de execução externa (estrutura na memória RAM que armazena endereços de retorno de sub-rotinas sem se limitar a registradores internos fixos), viabilizando o computador Altair 8800 em 1975 \cite{noyce1981}. Em 1976, já na Zilog, Faggin e Shima lançaram o Z80 com circuito de refrescamento para DRAM (lógica embutida que renova periodicamente a carga das células capacitivas da memória dinâmica para evitar perda de dados) e alimentação unipolar (operação com fonte única de \SI{+5}{\volt}, dispensando trilhos negativos adicionais), tornando o Z80 a base para sistemas com o CP/M, sistema operacional padrão dos primeiros microcomputadores de 8 bits antes do MS-DOS \cite{chmz80panel2007}. Em contraste, o processador MOS 6502, criado em 1975 por Chuck Peddle a partir do Motorola 6800, seguiu um ramo arquitetural independente, ilustrando uma trajetória paralela na expansão da computação pessoal, sem vínculo direto com a linhagem SGT iniciada no 4004.

\begin{figure}[!t]
\centering
\includegraphics[width=\columnwidth]{linhagem_faggin_cpu.pdf}
\caption{Linhagem arquitetural e de processo do Intel 4004 ao Zilog Z80 (1971 a 1976), destacando o MOS 6502 (1975) como ramo independente derivado do Motorola 6800.}
\label{fig:linhagem_faggin_cpu}
\end{figure}
